\subsection{缺陷试样评估}
\label{sec:results-defect}

\begin{figure}[tb]
\centering
\includegraphics[width=0.9\linewidth]{fig-bscan-snr-comparison.pdf}
\caption{16次平均下的B扫描成像结果：(a) PMDF-Net，(b) \placeholder{method-b}，(c) \placeholder{method-c}，(d) 常规延迟叠加。PMDF-Net产生最清晰的缺陷边界和最高对比度。缺陷位置由箭头指示。}
\label{fig:bscan-defect}
\end{figure}

\textbf{B扫描成像质量。}图~\ref{fig:bscan-defect}展示了在机加工缺陷试样上以16次平均（协议中最低堆叠水平）获取的代表性B扫描图像。
PMDF-Net对全部五个机加工槽产生最清晰的表征，缺陷边缘在各A扫描间锋利且空间连续。
相比之下，两种基线方法产生的边界明显更模糊、边缘定义更低，而常规延迟叠加波束形成器呈现最高背景散斑和最破碎的缺陷轮廓。
PMDF-Net的定性优越性在全部五个缺陷目标和全部测试的平均水平范围内保持一致。

\textbf{缺陷几何恢复。}机加工槽横截面尺寸为\SI{1.0}{\mm}$\times$\SI{20}{\mm}；PMDF-Net正确恢复了每个槽的侧向范围和深度穿透，未产生系统性的延长或缩短。远侧壁切口（切割至\SI{1.5}{\mm}深度）在PMDF-Net输出中清晰可辨，但在相同16次平均采集下两种基线方法均无法可靠检测。
这一发现表明PMDF-Net保留了与锐利几何特征相关的高波数区域空间频率内容——而这正是常规相干加权复合中首先退化的部分。

\textbf{定量指标。}表~\ref{tab:defect-metrics}总结了在三个成熟指标上的成像性能：对比度信噪比(CNR)、缺陷信噪比(SNR)和自动可检测性F1分数。PMDF-Net实现最高CNR为\placeholder{XX.XX}dB，优于最佳基线\placeholder{YY}dB。缺陷SNR遵循相同排序，PMDF-Net为\placeholder{ZZ.Z}dB versus\placeholder{W.W}dB for the next best method。五槽目标集的F1分数达到\placeholder{0.XXX}，而两种基线分别为\placeholder{0.YYY}和\placeholder{0.ZZZ}。三项指标均以显著差距倾向PMDF-Net。

\begin{table}[tb]
\centering
\caption{16次平均下机加工缺陷试样的成像性能指标。值为五个试样的均值$\pm$标准差。}
\label{tab:defect-metrics}
\begin{tabular}{lccc}
\toprule
方法 & CNR (dB) & 缺陷SNR (dB) & F1分数 \\
\midrule
PMDF-Net & \placeholder{XX.XX} $\pm$ \placeholder{s.x} & \placeholder{ZZ.Z} $\pm$ \placeholder{s.z} & \placeholder{0.XXX} $\pm$ \placeholder{s.f} \\
\placeholder{method-b} & \placeholder{WW.WW} $\pm$ \placeholder{s.w} & \placeholder{YY.Y} $\pm$ \placeholder{s.y} & \placeholder{0.YYY} $\pm$ \placeholder{s.g} \\
\placeholder{method-c} & \placeholder{VV.VV} $\pm$ \placeholder{s.v} & \placeholder{XX.X} $\pm$ \placeholder{s.x} & \placeholder{0.ZZZ} $\pm$ \placeholder{s.h} \\
延迟叠加 & \placeholder{UU.UU} $\pm$ \placeholder{s.u} & \placeholder{WW.W} $\pm$ \placeholder{s.w} & \placeholder{0.WWW} $\pm$ \placeholder{s.i} \\
\bottomrule
\end{tabular}
\end{table}

\textbf{统计显著性。}配合Holm-Bonferroni校正的成对Wilcoxon符号秩检验确认，PMDF-Net在全部三项指标上以$p < 0.05$水平显著优于每种基线方法。效应量（计算为等级-二列相关系数$r$）在每次比较中均超过0.5，满足大效应阈值。这些结果在五个试样上各五次重复扫描中保持一致。

\textbf{泛化发现。}PMDF-Net完全在完整试样上训练，未在缺陷数据集上做任何微调即进行评估。它超越的基线方法有的在同一完整数据上训练，
有的在完整-缺陷混合语料上训练——这一事实确认，多域物理正则化防止了对训练分布的过拟合，并实现向几何不同评估目标的可迁移。
这种零样本泛化是所提方法的首要实际优势。
